\section{Construction of the triangular power identities}
\label{sec:assembly}

We now construct one new triangular identity while preserving the low-degree
relations from Section~\ref{sec:low-relations}.  The division theorem allows
all the binary addition matrices to be used in the same coefficient algebra.
The remaining tasks are to make the final output invariant and to keep the
formal relation algebra finite over the old one.

We first record a binary addition using only linear and quadratic
relations.  The resulting algebra is finite over the input algebra;
Lemma~\ref{lem:sparsity-preserved} will then preserve the condition that
the characteristic set contains no curve of small degree.

\begin{lemma}[Finite extension from the Veronese relations]
\label{lem:veronese-extension}
Let \(b\ge2\), let \(k\) contain a primitive \(b\)-th root of unity
\(\zeta\), and let \(A\) be a commutative \(k\)-algebra with elements \(u,v\).
Introduce symbols \(T_\alpha\), indexed by
\(\alpha\in\Z_{\ge0}^{b}\) with \(|\alpha|=b\), and put
\[
 D_j=T_{b e_j}\quad(0\le j<b),\qquad
 z=\frac1b\sum_{j=0}^{b-1}D_j,\qquad
 u_T=-\frac1b\sum_{j=0}^{b-1}\zeta^{-j}D_j,\qquad
 v_T=T_{(1,\ldots,1)}.
\]
Let \(A'\) be the quotient of \(A[T_\alpha]\) by the relations
\begin{align}
 T_\alpha T_\beta&=T_\gamma T_\delta
       &&\text{whenever }\alpha+\beta=\gamma+\delta,
       \label{eq:veronese-exchanges}\\
 D_j&=z-\zeta^j u_T
       &&(0\le j<b),\qquad u_T=u,\quad v_T=v.
       \label{eq:veronese-linear}
\end{align}
Then \(A'\) is finite over \(A\), and in \(A'\)
\begin{equation}
 z^b=u^b+v^b,\qquad
 T_\alpha^b=\prod_{j=0}^{b-1}D_j^{\alpha_j}.
 \label{eq:veronese-monic}
\end{equation}
If \(A\) is graded and \(u,v\) have degree one, these are homogeneous
relations when each \(T_\alpha\) has degree one; the defining relations
\eqref{eq:veronese-exchanges}--\eqref{eq:veronese-linear} have degrees two
and one.
\end{lemma}

\begin{proof}
The pair exchanges identify any two products of the same number of
symbols whose total exponent vectors agree.  To see this directly,
regard such a product as an allocation of exponent units among factors,
each of total degree \(b\).  Choose the exponent vector desired in the
first factor.  Whenever that factor has an excess in one coordinate and
a deficit in another, some other factor has a unit in the deficient
coordinate.  Exchange this unit for a unit in the excess coordinate.
This is one relation of the form \eqref{eq:veronese-exchanges}.
Once the first factor has its desired vector, leave it fixed and continue
with the remaining factors.

Applying this observation to \(T_\alpha^b\) proves the second identity in
\eqref{eq:veronese-monic}.  In particular,
\[
 v^b=v_T^b=\prod_{j=0}^{b-1}D_j
       =\prod_{j=0}^{b-1}(z-\zeta^j u)=z^b-u^b.
\]
Thus \(z\) is integral over \(A\).  Every \(D_j\) belongs to \(A[z]\), and
the second identity in \eqref{eq:veronese-monic} makes every \(T_\alpha\)
integral over \(A[z]\).  There are finitely many symbols, so \(A'\) is
finite over \(A\).  The grading assertion follows from the displayed
defining relations.
\end{proof}

\begin{proposition}[The induction step]
\label{prop:induction-step}
Fix \(1\le i\le n\).  Suppose that a field \(k\), a central division algebra
\(\Delta_0\) over \(k\), a formal space \(U\), and a low-relation ideal
\(Q\subseteq\Sym_k U\) satisfy the induction conditions of
Section~\ref{sec:low-relations}, with the triangular identities already
constructed for the rows \(d>i\).
Then there are a finitely generated extension \(F/k\), a central division
algebra \(\Delta\) over \(F\), and an enlarged formal space and low-relation
ideal over \(F\) satisfying the same conditions with the rows \(d\ge i\)
constructed.  In particular, the enlarged formal space contains a form
\(t_i\) for which
\begin{equation}
 t_i^{a_i}
 =c_i f_i+\sum_{l<i}B_{il}f_l+\sum_{d>i}C_{id}t_d
 \pmod{Q_{\mathrm{new}}},
 \qquad c_i\in F^\times.
 \label{eq:new-triangular-row}
\end{equation}
All the enlarged formal forms evaluate to mutually commuting linear forms
in \(\Delta[x]\), and the new relation ideal is generated by linear forms
and quadrics.
\end{proposition}

\begin{proof}
Put \(b=a_i\), let \(e_0\) be the degree of \(\Delta_0\), and take
\(\mathfrak j\) and \(W\) as in
\eqref{eq:stage-ideal} and \eqref{eq:stage-annihilator}.
Write \(s=\dim_k W\).
Lemma~\ref{lem:stage-basepoint} gives base-point-freeness of
\(\mathfrak j_b\), and Lemma~\ref{lem:simple-factor} shows that every
nonzero geometric member of \(W\) has a factor of multiplicity one.

\smallskip\noindent
\emph{Parameters and the invariant sum.}
Choose integers in the following order:
\begin{align}
 h&\ge
 \max\left(\{2\}\cup
   \{v_p(e_0)s(p^{v_p(b)}-1):p\mid e_0\}\right),\label{eq:assembly-h}\\
 M&\ge
 \max\left\{2h+4s+10,\ \dim_k\Sym^b_k U\right\}.
 \label{eq:assembly-M}
\end{align}
The prime-indexed set in \eqref{eq:assembly-h} is empty if \(e_0=1\).
Theorem~\ref{thm:parameters} supplies the geometrically integral
projective constraint variety
\[
 X=\PP\left\{(\ell_1,\ldots,\ell_M)\in U^{\oplus M}:
           \sum_{m=1}^M w(\ell_m)=0\ \text{for every }w\in W\right\}.
\]
Use the curve and its two characters \(\chi,\psi\) from
Proposition~\ref{prop:binary-curves}.  For \(M-1\) copies of the curve,
indexed by \(m=2,\ldots,M\), let
\[
 H=G^{M-1},\qquad e_1=|H|,
\]
and write \(\chi_m,\psi_m:H\to\mu_b\) for the corresponding characters.
Define characters recursively by
\begin{equation}
 \lambda_M=1,\qquad
 \lambda_{m-1}=\lambda_m\chi_m\quad(M\ge m\ge2),\qquad
 \eta_1=\lambda_1,\quad \eta_m=\lambda_m\psi_m\quad(m\ge2).
 \label{eq:assembly-weights}
\end{equation}
Let \(H\) act on block \(m\) of the constraint cone by the point character
\(\eta_m^{-1}\).
This preserves its equations because their degree is \(b\), and it gives
the natural linearization of \(\cO_X(1)\).

As in Theorem~\ref{thm:division}, put
\[
 B=\prod_{m=2}^M\Pic^0(C),\qquad
 K=k(B\times X\times R),\qquad F=K^H,
\]
where \(R\) is an affine space with a faithful linear \(H\)-action.
Then \(K/F\) is Galois with group \(H\).
Let \(E=\bigotimes_{m=2}^M E_{A_m}\), of dimension \(e_1\) over \(K\),
and let \(\Delta_1\) be the descent of \(\End_K(E)\).
Theorem~\ref{thm:division}, applied with the block characters just chosen,
makes
\begin{equation}
 \Delta=(\Delta_0\otimes_k F)\otimes_F\Delta_1
 \label{eq:assembly-division}
\end{equation}
a division algebra.  Its scalar extension to \(K\) is
\(\Delta_0\otimes_k\End_K(E)\), and \(H\) fixes \(\Delta_0\) pointwise.
Notice that \(e_0,s,b\) were fixed before \(h\), and \(M\) before \(e_1\);
there is no restriction involving the new degree \(e_1\) in
\eqref{eq:assembly-h}.

The generic fiber of \(\cO_X(1)\), pulled back to \(K\), is a
one-dimensional semilinear \(H\)-space.  By Galois descent it has a nonzero
invariant generator \(\tau\).
Choose a \(k\)-basis of \(U\).  In each cone block, divide its coordinate
sections by \(\tau\) and use the resulting coefficients to form
\(\ell_m\in U_K\).
Transport of a coordinate section is pullback by \(g^{-1}\).
The inverse point action \(\eta_m^{-1}\) therefore gives
\begin{equation}
 g\ell_m=\eta_m(g)\ell_m.
 \label{eq:leaf-weights}
\end{equation}
In particular, dividing by the invariant \(\tau\) does not reverse the
character.

The constraint equations say that
\[
 S_b:=\sum_{m=1}^M\ell_m^b\in(\mathfrak j_K)_b.
\]
Since every \(\eta_m^b=1\), the element \(S_b\) is invariant and belongs to
\((\mathfrak j_F)_b\).
We claim that it has a decomposition
\begin{equation}
 S_b=c_i f_i+\sum_{l<i}B_{il}f_l+\sum_{d>i}C_{id}t_d
          \pmod{Q_F},\qquad c_i\in F^\times.
 \label{eq:invariant-sum}
\end{equation}
For this, let
\[
 V_b=Q_b+
       \sum_{l<i}f_l\Sym^{b-a_l}_k U+
       \sum_{d>i}t_d\Sym^{b-1}_k U.
\]
Then \(\mathfrak j_b=V_b+kf_i\).
If \(f_i\in V_b\), any prescribed nonzero coefficient of \(f_i\) may be
absorbed by changing the term in \(V_b\), proving the claim.

Suppose instead that \(f_i\notin V_b\).
Pure \(b\)-th powers span \(\Sym^b U\) in characteristic zero.
Over an algebraic closure, express \(f_i\) as a linear combination of at
most \(\dim\Sym^b U\) such powers, absorb each scalar into a \(b\)-th root,
and pad with zero forms.  The bound \eqref{eq:assembly-M} gives a nonzero
cone point with
\(\sum_m\ell_m^b=f_i\).
Thus the homogeneous coefficient of the class of \(f_i\) in
\(\mathfrak j_b/V_b\) is not identically zero on \(X\).
Geometric integrality makes it nonzero at the generic point.
Dividing all the coordinates by \(\tau\) only multiplies that coefficient
by a nonzero scalar to the \(b\)-th power.
Finally, the spaces \(V_b\) and \(\mathfrak j_b\) are defined over \(k\).
Invariant linear algebra over \(F\) gives
\eqref{eq:invariant-sum}, with homogeneous multipliers of the indicated
degrees.

\smallskip\noindent
\emph{The chain of binary additions over \(K\).}
We next enlarge the formal space over \(K\), keeping the old equations.
Start with \(s'_1=\ell_1\).
In link \(m\), introduce a fresh symbol \(T_{m,\alpha}\) for every
degree-\(b\) monomial \(p^\alpha\) in the curve coordinates.
Each symbol has formal degree one.
Evaluate it by the multiplication matrix from
Proposition~\ref{prop:binary-curves}, on the \(m\)-th tensor factor of \(E\),
with substitution
\begin{equation}
 U_1=s'_{m-1},\qquad U_2=\ell_m.
 \label{eq:link-substitution}
\end{equation}
Let \(s'_m\) be the formal linear combination of the diagonal symbols
corresponding to \(z\).

These substitutions preserve commutativity of the entire list.
Indeed, the coefficient matrices of a new link commute coefficientwise
with \(\Delta_0\) and with the earlier tensor factors.
The two inputs in \eqref{eq:link-substitution} commute with one another
and with all earlier forms.
The inputs therefore centralize the new coefficient algebra
\(\End_K(E_{A_m})\), embedded in the \(m\)-th tensor factor.
Consequently substitution defines a \(K\)-algebra homomorphism
\[
 \begin{aligned}
 \End_K(E_{A_m})[U_1,U_2]
 &\longrightarrow
 \bigl(\Delta_0\otimes_k\End_K(E)\bigr)[x_1,\ldots,x_n],\\
 U_1&\longmapsto s'_{m-1},\qquad U_2\longmapsto\ell_m,
 \end{aligned}
\]
where the inputs on the right denote their evaluated forms.
It preserves the polynomial commutation identities of the new
multiplication matrices, and every new evaluated form commutes with
every earlier one.
This argument uses no coefficientwise commutativity among the matrices
within a single link.

For each link impose the pair-exchange quadrics
\eqref{eq:veronese-exchanges}, the diagonal relations
\eqref{eq:veronese-linear}, and the identifications
\[
 u_T=s'_{m-1},\qquad v_T=\ell_m.
\]
They are identities under evaluation because they are identities of
sections on the curve.
Lemma~\ref{lem:veronese-extension} shows that the resulting quotient is
finite over the preceding formal relation algebra and gives
\[
 (s'_m)^b=(s'_{m-1})^b+\ell_m^b.
\]
After all links, the formal quotient is finite over
\(\Sym_K U_K/Q_K\), is defined by the old and new linear and quadratic
relations, and satisfies
\begin{equation}
 (s'_M)^b=\sum_{m=1}^M\ell_m^b=S_b.
 \label{eq:telescoping-links}
\end{equation}

\smallskip\noindent
\emph{Equivariance and descent.}
Give the new symbol space in link \(m\) the degree-\(b\) action on the
\(p\)-monomials from Proposition~\ref{prop:binary-curves}, multiplied by
the character \(\lambda_m\), and extend it semilinearly over \(K\).
Keep \(U\) fixed.
The pair-exchange relations are stable under the monomial phase and
permutation action and under the common twist.
The diagonal relations are stable as well.
The distinguished symbols \(u_T,v_T,z\) have respective weights
\[
 \lambda_m\chi_m=\lambda_{m-1},\qquad
 \lambda_m\psi_m=\eta_m,\qquad
 \lambda_m.
\]
These are the weights of their prescribed inputs and output:
the input identifications are stable by induction, starting from
\eqref{eq:leaf-weights}.  Thus the full ideal of linear and quadratic
relations is stable.

Evaluation is equivariant, not merely the relation ideal.
For a section \(a\) before the twist, write
\(M_a(U_1,U_2)\) for its linear multiplication matrix.
The transport identity \eqref{eq:matrix-transport} reads
\[
 M_{g_*a}(U_1,U_2)
 =(g_{\mathrm{coeff}}M_a)
      \bigl(\chi_m(g)U_1,\psi_m(g)U_2\bigr),
\]
where \(g_{\mathrm{coeff}}\) includes the field action and transport on
the corresponding endomorphism factor.
Scalar choices in transporting the universal line cancel on these
endomorphisms.
Using the weights of the two inputs and the degree-one homogeneity of
the matrix polynomial, we obtain
\begin{align*}
 g\!\left(M_a(s'_{m-1},\ell_m)\right)
 &=(g_{\mathrm{coeff}}M_a)
       \bigl(\lambda_m(g)\chi_m(g)s'_{m-1},
             \lambda_m(g)\psi_m(g)\ell_m\bigr)\\
 &=\lambda_m(g)\,M_{g_*a}(s'_{m-1},\ell_m).
\end{align*}
This is exactly the action assigned to the twisted symbol.
It proves equivariance successively through all the links.

Take invariant vector spaces under semilinear Galois descent.
The enlarged formal space descends to an \(F\)-space containing \(U_F\),
and the full degree-one and degree-two pieces of the stable relation
ideal descend to finite-dimensional spaces over \(F\). Let
\(Q_{\mathrm{new}}\) be the ideal generated by these descended spaces.
Scalar extension commutes with the multiplication maps generating this
ideal. Since the original relation ideal was generated in degrees one
and two, extension to \(K\) recovers exactly that ideal.
Equivariant evaluation takes these invariant forms into
\(\Delta[x]\), with \(\Delta\) as in \eqref{eq:assembly-division}.
They are mutually commuting.

The last output \(s'_M\) is invariant because \(\lambda_M=1\); define
\(t_i=s'_M\).
Combining \eqref{eq:invariant-sum} with
\eqref{eq:telescoping-links} proves \eqref{eq:new-triangular-row}.
All earlier equations are retained, so every old triangular identity
continues to hold modulo \(Q_{\mathrm{new}}\).

It remains to check the other induction conditions.
Finiteness of the new relation algebra over the old one can be checked
after the faithfully flat extension \(F\subset K\), where it was proved
link by link.  In particular it remains finite over \(F[x]\).
Lemma~\ref{lem:sparsity-preserved} now preserves the geometric
no-small-curve condition: the linear projection of the new
characteristic set is finite and pulls back the old hyperplane bundle.
The central variables still evaluate to themselves, and the formal space
still contains \(U_F\).
Finally \(K\) is finitely generated over \(k\), and so is its finite-group
invariant field \(F\).
Thus the center is still finitely generated over the original
\(\C\), and all induction conditions have been verified.
\end{proof}

\begin{proof}[Proof of Theorem~\ref{thm:construction}]
Proposition~\ref{prop:initial-data} supplies the initial data with no
triangular rows yet constructed.
Apply Proposition~\ref{prop:induction-step} in the order
\(i=n,n-1,\ldots,1\).
Every step retains all preceding identities as consequences of an ideal
generated in degrees at most two.  Thus a later, smaller value of \(a_i\)
does not lose an earlier power identity.
After finitely many steps we have the division algebra and commuting
linear forms satisfying \eqref{eq:triangular} for all \(i\).
Let \(A\) be the commutative graded algebra generated by all the evaluated
formal forms. The center is
finitely generated over \(\C\), all multipliers in the triangular
identities belong to \(A\) in the stated degrees, and all \(c_i\) are
nonzero. Lemma~\ref{lem:ordered-basis} gives finiteness over \(k[x]\)
and the full ordered left-\(\Delta\) basis. This proves every assertion
of Theorem~\ref{thm:construction}.
\end{proof}

\begin{proof}[Proof of Corollary~\ref{thm:main}]
Apply Theorem~\ref{thm:construction} to the regular sequence contained
in \(I\). Proposition~\ref{prop:monomial-reduction} gives one monomial
ideal \(J\) containing all \(x_i^{a_i}\) and having the same Hilbert
function as \(I\) in every degree.
\end{proof}
